\documentclass[11pt]{article}
% Adapted from an original template by Hlyni Arnórssyni,
% Reykjavik University, Iceland

% ------------------------------ PAGE SETTINGS ------------------------------
\usepackage{geometry}

\geometry{
    paper=a4paper,
    top=2.5cm,
    bottom=2.5cm,
    left=2.5cm,
    right=2.4cm,
    headheight=0.75cm,
    footskip=1.5cm,
    headsep=0.75cm
}

% ------------------------------ BASIC PACKAGES ------------------------------
\usepackage{blindtext}

% Character encoding
\usepackage[T1]{fontenc}
\usepackage[utf8]{inputenc}

% Mathematics
\usepackage{amsmath, amsfonts, amsthm, amssymb}
\usepackage{braket, nicefrac}

% SI units
\usepackage{siunitx}

% Lists / columns
\usepackage{enumitem, multicol}

% ------------------------------ FIGURES ------------------------------
\usepackage{graphicx}
\usepackage{float}
\usepackage{keystroke}
\usepackage{pgfplots}
\usepgfplotslibrary{units}
\pgfplotsset{compat=1.16}

% Search these locations automatically for figures
\graphicspath{{fig/}{graphics/}{Graphics/}{./}}

% ------------------------------ TABLES ------------------------------
\usepackage{tabularx}
\usepackage{array}
\usepackage[table]{xcolor}

% ------------------------------ TYPOGRAPHY ------------------------------
\usepackage{caption}
\usepackage{titlesec}
\usepackage{microtype}
\usepackage{csquotes}

% ------------------------------ COLORS ------------------------------
\definecolor{GWBlue}{HTML}{1E4E79}

% ------------------------------ SECTION STYLING ------------------------------
\titleformat{\section}
  {\Large\bfseries\color{GWBlue}}
  {\thesection.}
  {0.40em}
  {}

\titlespacing*{\section}
  {0pt}
  {0.70em}
  {0.28em}

% ------------------------------ CAPTION STYLING ------------------------------
\captionsetup{
    font=small,
    labelfont=bf,
    textfont=it,
    labelsep=period,
    skip=4pt
}

% ------------------------------ PARAGRAPH SPACING ------------------------------
\setlength{\parindent}{0pt}
\setlength{\parskip}{0.30em}

% ------------------------------ BIBLIOGRAPHY ------------------------------
\usepackage[
    backend=biber,
    style=authoryear,
    sorting=none,
    giveninits=true,
    maxbibnames=99,
    dashed=false
]{biblatex}

\addbibresource{references.bib}

\DeclareNameAlias{sortname}{family-given}
\DeclareNameAlias{default}{family-given}
\DeclareFieldFormat[article]{title}{#1}
\DeclareFieldFormat[misc]{title}{#1}
\renewbibmacro{in:}{}

\AtEveryBibitem{%
    \clearfield{doi}%
    \clearfield{urldate}%
}

\AtBeginBibliography{\small}
\setlength{\bibitemsep}{0.35em}

% ------------------------------ HYPERLINKS ------------------------------
\usepackage{hyperref}

\hypersetup{
    colorlinks=true,
    linkcolor=GWBlue,
    citecolor=GWBlue,
    urlcolor=GWBlue
}
\usepackage{booktabs}
\newcommand{\optfig}[2]{\IfFileExists{#2}{\includegraphics[#1]{#2}}{\fbox{\parbox{0.9\linewidth}{\centering\small Figure not copied yet [TBD]:\\\scriptsize\nolinkurl{#2}}}}}

\begin{document}

\begin{titlepage}
\centering
\thispagestyle{empty}

\vspace*{0.15in}

\includegraphics[
    width=0.5\textwidth
]{gwu_logo.pdf}

\vspace{0.28in}

{\large\bfseries Data Science Program\par}
{\large\bfseries Department of Data Science\par}
{\large DATS 6501 - Capstone Project - Fall 2026\par}

\vspace{0.38in}

{\Large\bfseries
Post-Training Analysis \& Novelty Detection Using\par}

\vspace{0.08in}

{\Large\bfseries
Neural Network Self-Organizing Maps (NNSOM)\par}

\vspace{0.34in}

{\large\bfseries
PHASE 3 REPORT\par}

\vspace{0.22in}

{\large\bfseries
Novelty Detection\par}
\vspace{0.12in}
{\large\bfseries Fashion-MNIST\par}

\vspace{0.36in}

{\large\bfseries Capstone - Group 5\par}

\vspace{0.08in}

{\large\bfseries Fyrooz Anika Khan\par}
{\large\bfseries Nazish Atta\par}

\vspace{0.30in}

{\large The George Washington University\par}
{\large Washington, DC\par}

\vspace{0.30in}

{\large\itshape Supervised by\par}
{\large\itshape Dr. Amir Jafari\par}
{\large\itshape Associate Professor of Data Science\par}

\vspace*{\fill}

\end{titlepage}

\setcounter{page}{2}
\newcommand{\tbd}{\textbf{[TBD]}}
\setlength{\emergencystretch}{3em}

\section{Overview}

Phase 3 asks whether the frozen Fashion-MNIST network and its 15$\times$15 SOM can tell when an image is \emph{unlike} the training data (a ``novel'' or out-of-distribution image). The design is the MNIST Phase 3 design, copied unchanged and written down before any Fashion-MNIST novelty score was computed. There are two tests:
\begin{itemize}[leftmargin=1.4em, itemsep=1pt]
  \item \textbf{Far novelty:} the 26,032 SVHN test images (house-number photographs \cite{netzer2011reading}) against the 10,500 Fashion-MNIST validation images. These are the same SVHN images used for MNIST.
  \item \textbf{Near novelty:} clothing classes the network has never seen. In 5 class splits a new network and map are trained on 6 classes and tested on the other 4 \cite{neal2018open}.
\end{itemize}
The Fashion-MNIST test split is not used in Phase 3.

\textbf{Summary of the results.}
\begin{itemize}[leftmargin=1.4em, itemsep=1pt]
  \item \textbf{The SOM detects SVHN, but weakly.} Its quantization error (QE) gives AUROC 0.789, well above chance (0.500), but far below the network's energy score (0.982) and max-softmax (0.884). On MNIST, QE gave 0.981 on the same SVHN images.
  \item \textbf{The fixed QE alarm does not work on SVHN.} At the threshold set from the training images it flags only 3.0\% of SVHN, fewer than the 5.1\% of normal validation images it flags. On MNIST it caught 96.2\%.
  \item \textbf{SVHN is mostly told apart by one simple feature.} Its embeddings are much shorter than Fashion-MNIST's; length alone separates the two sets almost perfectly (AUROC 0.995). QE follows length only weakly, which explains its weak result.
  \item \textbf{On unseen clothing classes the SOM beats max-softmax.} QE is higher than max-softmax in all 5 splits (+0.046 AUROC, 95\% interval +0.034 to +0.057). On MNIST this comparison was inconclusive.
  \item \textbf{Plain nearest-neighbour distance beats the SOM in both tests.}
\end{itemize}

\subsection*{Where each item of the plan is covered}

{\small
\renewcommand{\arraystretch}{1.12}
\begin{tabularx}{\textwidth}{|>{\raggedright\arraybackslash}p{0.52\textwidth}|X|}
\hline
\textbf{What Phase 3 asks for} & \textbf{Where} \\ \hline
QE for every image; threshold at the 95th percentile of training QE & Sections~\ref{sec:method} and~\ref{sec:qeflag}, Table~\ref{tab:sweep} \\ \hline
U-matrix as a novelty signal; novelty scores shown on the SOM grid & Section~\ref{sec:umatrix}, Figures~\ref{fig:valqe}--\ref{fig:missed} \\ \hline
Compare with Isolation Forest, One-Class SVM and an autoencoder & Section~\ref{sec:baselines}, Table~\ref{tab:classical} \\ \hline
Compare with max-softmax, energy, Mahalanobis (same embeddings, AUROC) & Section~\ref{sec:far}, Table~\ref{tab:far} \\ \hline
AUROC, AUPR, FPR at 95\% TPR & Table~\ref{tab:far} \\ \hline
Failure modes: which novel images the SOM misses, and why & Sections~\ref{sec:umatrix} and~\ref{sec:length} \\ \hline
\end{tabularx}}

\section{Key terms in plain words}

{\small
\renewcommand{\arraystretch}{1.12}
\begin{tabularx}{\textwidth}{|>{\raggedright\arraybackslash}p{0.26\textwidth}|X|}
\hline
\textbf{Term} & \textbf{What it means} \\ \hline
Normal / novel images & Normal images come from the same source as the training data (Fashion-MNIST validation). Novel images do not (SVHN, or clothing classes the network never saw). \\ \hline
Far / near novelty & Far: novel images that look very different (street-number photos). Near: novel images that look similar (another kind of clothing). \\ \hline
Quantization error (QE) & The distance from an image's embedding to its best SOM cell. A large QE means the image does not fit the map well. \\ \hline
U-matrix score & How far the image's cell is from its neighbouring cells; high on the borders between groups. \\ \hline
Max-softmax, energy & Two scores from the network's own output. Low confidence (max-softmax) or high energy suggests a novel image \cite{hendrycks2017baseline,liu2020energy}. \\ \hline
Mahalanobis, nearest neighbour & Two distances in the embedding space: to the class centres \cite{lee2018mahalanobis}, and to the 5 closest training embeddings \cite{sun2022knn}. \\ \hline
AUROC & How well a score ranks novel images above normal ones: 1 is perfect, 0.5 is chance, below 0.5 means the score is the wrong way round. \\ \hline
AUPR-out & Precision–recall area with novel images as positives. Its chance level here is 0.713, because 71\% of the images are SVHN. \\ \hline
FPR at 95\% TPR & The share of normal images wrongly flagged when 95\% of the novel images are caught. Lower is better. \\ \hline
Paired interval & A 95\% range for the difference between two scores, computed on the same resampled images. \\ \hline
\end{tabularx}}

\section{Related Work}

Kohonen's SOM \cite{kohonen1990som} gives a set of prototypes on a grid, and the distance to the nearest prototype (QE) is a natural novelty signal; we use the NNSOM implementation \cite{jafari_nnsom} and do not change the map. Hendrycks and Gimpel \cite{hendrycks2017baseline} showed that the network's top softmax probability is a strong, simple novelty baseline, and Liu et al.\ \cite{liu2020energy} replaced it with an energy score computed from all the logits. Lee et al.\ \cite{lee2018mahalanobis} measure the distance to class centres in feature space, and Sun et al.\ \cite{sun2022knn} the distance to the nearest training features; we compute both on the same \texttt{fc2} embeddings as the SOM, so the comparison isolates what the map adds. Isolation Forest \cite{liu2008isolation}, the One-Class SVM \cite{scholkopf2001support} and autoencoder reconstruction error \cite{sakurada2014autoencoder} are the classical detectors named in the project proposal. Nalisnick et al.\ \cite{nalisnick2019deep} showed that a model of normal data can rate SVHN as \emph{more} typical than its own data; we see the same reversal for two classical detectors. For near novelty we follow the open-set protocol of Neal et al.\ \cite{neal2018open}, where max-softmax is a hard baseline to beat \cite{vaze2022openset}.

\section{Method}\label{sec:method}

\begin{itemize}[leftmargin=1.4em, itemsep=2pt]
\item \textbf{Frozen inputs.} The Phase 1 network (\texttt{lenet5\_fashion\_mnist\_s42}) and the Phase 2 headline SOM (\texttt{som\_15x15\_nb11\_ep250\_s42}), each checked by its SHA-256 before every run.
\item \textbf{Outside set.} All 26,032 SVHN test images, in file order, converted exactly as for MNIST: RGB to grey, resized from $32\times32$ to $28\times28$, pixels scaled to $[0,1]$. Their digit labels are not used.
\item \textbf{Scores.} S1 QE, S2 U-matrix, S3 max-softmax, S4 energy, S5 Mahalanobis (shared Ledoit--Wolf covariance \cite{ledoit2004wellconditioned}), S6 mean distance to the 5 nearest training embeddings. Anything fitted is fitted on the training embeddings only.
\item \textbf{Measures.} AUROC with a bootstrap 95\% interval (2,000 resamples, seed 42) \cite{fawcett2006roc}, AUPR-out \cite{saito2015pr}, FPR at 95\% TPR, and paired bootstrap intervals for QE minus each other score. Control: the QE AUROC with the labels shuffled 1,000 times.
\item \textbf{Statements, fixed before the run.} O1: the SOM detects SVHN (QE interval above 0.5). \textbf{O2 (primary):} the SOM is competitive with max-softmax (lower end of the QE minus max-softmax interval above $-0.01$). O3: the same against energy, Mahalanobis and nearest neighbour, one by one. O4: the U-matrix detects SVHN. Each is met, not met, or inconclusive.
\item \textbf{QE alarm.} An image is flagged if its QE is above the 95th percentile of the training QE (1.4215). Other percentiles are reported only as a sweep; none is chosen from it.
\end{itemize}

\section{Far Novelty: SVHN}\label{sec:far}

\begin{table}[H]
\centering\footnotesize\setlength{\tabcolsep}{4pt}
\begin{tabular}{@{}lcccc@{}}\toprule
Score & AUROC [95\% interval] & AUPR-out & FPR at 95\% TPR & QE minus score [95\% interval]\\\midrule
S4 Energy & 0.982 [0.981, 0.984] & 0.988 & 0.063 & $-0.193$ [$-0.199$, $-0.186$]\\
S3 Max-softmax & 0.884 [0.880, 0.888] & 0.936 & 0.307 & $-0.094$ [$-0.101$, $-0.088$]\\
S6 Nearest neighbour & 0.864 [0.859, 0.869] & 0.898 & 0.314 & $-0.074$ [$-0.078$, $-0.070$]\\
S2 U-matrix & 0.812 [0.805, 0.818] & 0.892 & 0.253 & $-0.022$ [$-0.030$, $-0.014$]\\
S1 SOM QE & 0.789 [0.783, 0.796] & 0.825 & 0.420 & --\\
S5 Mahalanobis & 0.713 [0.706, 0.720] & 0.768 & 0.454 & $+0.076$ [$+0.071$, $+0.082$]\\\bottomrule
\end{tabular}
\caption{Fashion-MNIST validation (10,500, normal) against SVHN test (26,032, novel). One network and one map. AUPR-out chance level: 0.713. Run \texttt{fashion\_mnist\_svhn\_novelty\_statements\_s42\_v1}.}
\label{tab:far}
\end{table}

\begin{table}[H]
\centering\small
\begin{tabular}{@{}lll@{}}\toprule
Statement & Result & Verdict\\\midrule
O1: the SOM detects SVHN & QE AUROC 0.789, interval above 0.5 & \textbf{met}\\
\textbf{O2 (primary): competitive with max-softmax} & $-0.094$, upper end below $-0.01$ & \textbf{not met}\\
O3: competitive with energy & $-0.193$ & not met\\
O3: competitive with Mahalanobis & $+0.076$ & met\\
O3: competitive with nearest neighbour & $-0.074$ & not met\\
O4: the U-matrix detects SVHN & AUROC 0.812, interval above 0.5 & \textbf{met}\\\bottomrule
\end{tabular}
\caption{The statements, judged by the rules fixed before the run.}
\label{tab:statements}
\end{table}

\begin{itemize}[leftmargin=1.4em, itemsep=2pt]
\item \textbf{The SOM sees SVHN as unusual, but not strongly.} With shuffled labels the QE AUROC averages 0.500 (largest of 1,000 shuffles: 0.510), so 0.789 is far from chance (O1 met). The U-matrix score does a little better (0.812, O4 met).
\item \textbf{The network's own scores are clearly better} (O2 not met). Energy almost separates the two sets perfectly (0.982); max-softmax and nearest-neighbour distance are also ahead of QE. Only Mahalanobis is behind.
\item \textbf{Compared with MNIST.} On the same SVHN images the MNIST map's QE gave 0.981. So the SOM's far-novelty signal depends strongly on what the network was trained on.
\item \textbf{How stable is this?} The intervals above come from one map. In the Phase 2 feature-layer ablation, five SOMs trained on the same embeddings with different seeds gave QE AUROCs from 0.766 to 0.884. That spread is several times wider than the intervals, but QE stays well below energy at every seed. The same ablation found that maps on the earlier \texttt{fc1} layer detect SVHN better (QE AUROC 0.902 on average, against 0.815 for \texttt{fc2}), because \texttt{fc1} still carries more about how an image looks; \texttt{fc2} stays the headline layer because it gives the better map of the network's mistakes.
\end{itemize}

\section{The Fixed QE Alarm}\label{sec:qeflag}

At the 95th percentile of the training QE (1.4215) the alarm flags 536 of the 10,500 validation images (5.1\%) but only 792 of the 26,032 SVHN images (3.0\%). It flags \emph{fewer} novel images than normal ones, so as a yes/no detector it does not work here. Table~\ref{tab:sweep} shows that this holds for every stricter cut-off as well; only the looser 90th percentile flags slightly more SVHN than validation images. The SOM still ranks most SVHN images above most validation images (AUROC 0.789), but SVHN sits in the middle of the QE range, below the upper tail of the normal images, where any threshold must be placed.

\begin{table}[H]
\centering\small
\begin{tabular}{@{}ccccc@{}}\toprule
Training percentile & QE cut-off & Validation flagged & SVHN detected & SVHN missed\\\midrule
90 & 1.245 & 10.1\% & 13.1\% & 22,617\\
95 (fixed) & 1.421 & 5.1\% & 3.0\% & 25,240\\
97.5 & 1.599 & 2.7\% & 1.0\% & 25,760\\
99 & 1.846 & 1.1\% & 0.26\% & 25,965\\
99.5 & 2.023 & 0.6\% & 0.09\% & 26,009\\\bottomrule
\end{tabular}
\caption{QE cut-offs at different percentiles of the training QE. Descriptive; the 95th percentile was fixed in advance and no cut-off was chosen from this table. Run \texttt{fashion\_mnist\_svhn\_qe\_threshold\_sweep\_s42\_v1}.}
\label{tab:sweep}
\end{table}

\textbf{QE and the network's mistakes.} Within the validation images, the 536 flagged images contain 73 of the network's 876 mistakes (13.6\% of the flagged images are wrong, against 8.3\% overall). The same number of random picks would hold about 45 mistakes (5 random draws: 43 to 50; exact $p<0.001$). So a large QE does point to harder normal images, as on MNIST.

\begin{figure}[H]
\centering
\optfig{width=0.58\textwidth,keepaspectratio}{fig/qe_novelty_detection/qe_score_distribution.pdf}
\caption{Training and validation QE, with the fixed threshold at the 95th percentile of the training QE.}
\label{fig:qe}
\end{figure}

\section{U-Matrix and Where SVHN Lands}\label{sec:umatrix}

On the validation images, QE and the U-matrix score have a Spearman correlation of 0.38: they are related but measure different things. Figure~\ref{fig:overlay} shows where the SVHN images land on the map. All 26,032 fall in only 42 of the 225 cells, and two cells, 109 and 94, receive 11,977 and 11,298 images (89\% together). Both lie in the region whose most common Fashion-MNIST class is \emph{Bag}. In total 91\% of SVHN lands in Bag cells, and these hold 92\% of the missed images. A further 9\% lands in mixed Pullover and Dress cells. A missed SVHN image has a QE only slightly above that of the normal images in the same cell (median ratio 1.04). So, to the map, SVHN looks like a typical bag. This is a description of where the misses are, not a proof of why.

\begin{figure}[H]
\centering
\optfig{width=0.60\textwidth,keepaspectratio}{fig/u_matrix_analysis/fashion_mnist_validation_qe_som_overlay.pdf}
\caption{QE of the Fashion-MNIST validation images shown on the SOM grid (222 occupied cells). The missed SVHN images all have a QE at or below the threshold (median 1.07, range 0.50 to 1.42), inside the range of normal images. Run \texttt{fashion\_mnist\_svhn\_proposal\_gap\_analysis\_s42\_v1}.}
\label{fig:valqe}
\end{figure}

\begin{figure}[H]
\centering
\optfig{width=0.95\textwidth,keepaspectratio}{fig/u_matrix_analysis/svhn_som_overlay.pdf}
\caption{SVHN test images on the $15\times15$ SOM: images per cell, share missed at the fixed threshold, and the most common Fashion-MNIST class of each cell. Descriptive. Run \texttt{fashion\_mnist\_svhn\_som\_overlay\_s42\_v1}.}
\label{fig:overlay}
\end{figure}

\begin{figure}[H]
\centering
\begin{minipage}[c]{0.48\textwidth}\centering
\optfig{width=\linewidth,keepaspectratio}{fig/u_matrix_analysis/u_matrix_novelty_som_grid.pdf}
\end{minipage}\hfill
\begin{minipage}[c]{0.48\textwidth}\centering
\optfig{width=\linewidth,keepaspectratio}{fig/u_matrix_analysis/qe_u_matrix_validation.pdf}
\end{minipage}
\caption{Left: U-matrix of the $15\times15$ SOM. Right: QE against the U-matrix score on the validation images (Spearman 0.38).}
\label{fig:umatrix}
\end{figure}

\begin{figure}[H]
\centering
\optfig{width=0.70\textwidth,keepaspectratio}{fig/novelty_benchmark/qe_missed_svhn_examples.pdf}
\caption{SVHN images with the lowest QE, all missed at the fixed threshold. Descriptive.}
\label{fig:missed}
\end{figure}

\section{Why SVHN Is Missed: Embedding Length}\label{sec:length}

These checks were run because the same problem appeared as on MNIST, and are exploratory.
\begin{itemize}[leftmargin=1.4em, itemsep=2pt]
\item \textbf{SVHN embeddings are short.} The average length of an SVHN embedding is 8.5, against 20.7 for Fashion-MNIST; 95\% of SVHN images are below 11.7, and 95\% of validation images are above 13.2. A rule that scores ``shorter means more novel'' separates the two sets almost perfectly (AUROC 0.995).
\item \textbf{Not a brightness effect.} SVHN images are often dark digits on a bright background, the opposite of Fashion-MNIST. Inverting their grey levels lengthens the embeddings only from 8.50 to 9.18 on average.
\item \textbf{QE follows length only weakly.} With the validation and SVHN images pooled, QE and embedding length have a Spearman correlation of $-0.29$ (MNIST: $-0.83$); within the validation images it is $-0.02$ and within SVHN $+0.13$. Keeping only images whose length lies in the 5--95\% range of the training lengths (13.2 to 31.8) leaves 9,478 validation images but only 532 SVHN images. On that subset QE's AUROC is 0.958 (0.953--0.962), against 0.789 on all images. This fits the idea that the short embeddings are what hide SVHN from QE, but 532 selected images are too few to show it.
\end{itemize}
So on Fashion-MNIST, SVHN differs from the normal images mainly in one simple feature, which a trivial rule already uses. A far-novelty AUROC on SVHN largely measures how strongly a score follows that feature. On MNIST the length effect helped QE; on Fashion-MNIST it works against it.

\section{Classical Baselines}\label{sec:baselines}

\begin{table}[H]
\centering\footnotesize\setlength{\tabcolsep}{4pt}
\begin{tabular}{@{}lccccl@{}}\toprule
 & \multicolumn{3}{c}{One run (seed 42)} & Five seeds & \\
Method & AUROC & AUPR-out & FPR95 & AUROC (mean $\pm$ SD) & Note\\\midrule
One-Class SVM & 0.734 & 0.811 & 0.634 & $0.738 \pm 0.005$ & \\
Autoencoder, raw inputs & 0.583 & 0.702 & 0.744 & $0.395 \pm 0.030$ & \\
Autoencoder, standardised & 0.676 & 0.768 & 0.653 & $0.742 \pm 0.089$ & varies by seed \\
Isolation Forest & 0.266 & 0.564 & 0.934 & $0.272 \pm 0.036$ & reversed \\
Mahalanobis (refitted) & -- & -- & -- & $0.738 \pm 0.002$ & \\
Nearest neighbour (refitted) & -- & -- & -- & $0.859 \pm 0.011$ & \\
Random score (control) & -- & 0.714 & 0.948 & 0.505 & \\\bottomrule
\end{tabular}
\caption{AUROC of the classical detectors on the same images and embeddings. FPR95: false-positive rate at 95\% TPR; AUPR-out chance level 0.713. The five-seed column refits these methods on the same 10,000 training rows per seed (seeds 42--46); it is exploratory. The one-run autoencoder and One-Class SVM use all 49,000 and 10,000 training rows. Runs \texttt{fashion\_mnist\_svhn\_five\_method\_comparison\_s42\_v1}, \texttt{fashion\_mnist\_svhn\_ae\_train\_standardized\_s42\_v1}, \texttt{fashion\_mnist\_svhn\_fair\_budget\_baselines\_s42\_46\_v1} and \texttt{fashion\_mnist\_svhn\_unified\_nine\_methods\_s42\_46\_v1}.}
\label{tab:classical}
\end{table}

None of the classical detectors beats the SOM's QE by much, and two behave badly. \textbf{Isolation Forest} ranks SVHN as \emph{more normal} than Fashion-MNIST (AUROC 0.27), as it did on MNIST (0.24), and it stays below chance with standardised inputs and with 49,000 training rows. It treats the short SVHN embeddings as typical. \textbf{The autoencoder} depends mainly on input scaling: raw inputs give 0.40 to 0.55, standardised inputs 0.72 to 0.74, and single seeds range from 0.59 to 0.82. This matches the pattern reported for models of normal data on SVHN \cite{nalisnick2019deep}.

\section{Near Novelty: Unseen Clothing Classes}\label{sec:heldout}

For each of 5 class splits (the same class numbers as for MNIST, drawn at random once), a new LeNet-5 and a new $15\times15$ SOM are trained on 6 classes only; the validation images of the other 4 classes are the novel images. The networks reach 92.8\% to 96.6\% accuracy on their known classes. Statements: H1, the SOM detects the unseen classes; \textbf{H2 (primary), QE is better than max-softmax}; H3, QE is not worse than max-softmax, energy or Mahalanobis by more than 0.01; H4, QE is better than nearest-neighbour distance. Each uses a 95\% $t$-interval over the 5 per-split differences.

\begin{table}[H]
\centering\footnotesize\setlength{\tabcolsep}{3pt}
\begin{tabular}{@{}clcccccc@{}}\toprule
Split & Unknown classes & QE & MSP & Energy & Mahalanobis & $k$NN & Length\\\midrule
1 & 0, 1, 5, 6 & 0.801 & 0.750 & 0.781 & 0.828 & 0.829 & 0.675\\
2 & 2, 3, 5, 8 & 0.831 & 0.787 & 0.885 & 0.860 & 0.892 & 0.814\\
3 & 1, 5, 6, 7 & 0.739 & 0.695 & 0.678 & 0.814 & 0.812 & 0.634\\
4 & 2, 5, 6, 8 & 0.817 & 0.760 & 0.788 & 0.796 & 0.859 & 0.641\\
5 & 0, 4, 5, 8 & 0.783 & 0.751 & 0.782 & 0.814 & 0.819 & 0.670\\\midrule
\multicolumn{2}{c}{Mean} & 0.794 & 0.749 & 0.783 & 0.823 & 0.842 & 0.687\\
\multicolumn{2}{c}{QE minus score} & -- & $+0.046$ & $+0.011$ & $-0.028$ & $-0.048$ & --\\
\multicolumn{2}{c}{(95\% interval)} & & [$0.034$, $0.057$] & [$-0.041$, $0.064$] & [$-0.071$, $0.014$] & [$-0.071$, $-0.026$] & \\\bottomrule
\end{tabular}
\caption{Near novelty: AUROC for detecting the 4 unseen classes among the 6 known ones (validation images), one new network and map per split. Class numbers: 0 T-shirt/top, 1 Trouser, 2 Pullover, 3 Dress, 4 Coat, 5 Sandal, 6 Shirt, 7 Sneaker, 8 Bag, 9 Ankle boot. Within a split the bootstrap intervals are about $\pm0.008$ wide. Run \texttt{held\_out\_class\_novelty\_lenet5\_som15x15\_nb11\_ep250\_s42}.}
\label{tab:heldout}
\end{table}

\begin{itemize}[leftmargin=1.4em, itemsep=2pt]
\item \textbf{H1 shown.} QE AUROC 0.739 to 0.831, every interval above 0.5; the label-shuffle control gives 0.500 in every split.
\item \textbf{H2 shown: the SOM beats max-softmax.} QE is higher in all 5 splits, by 0.046 on average (0.034 to 0.057). On MNIST this was inconclusive. Clothing classes look alike, so the network is often confidently wrong about an unseen class, while the map still places it away from the known prototypes.
\item \textbf{Not a length effect here.} Length alone gives only 0.687 on average, below QE in every split.
\item \textbf{H3:} shown against max-softmax; inconclusive against energy and Mahalanobis (the splits vary too much).
\item \textbf{H4 shown in the opposite direction.} Nearest-neighbour distance beats QE in every split ($-0.048$). The 225 cells lose detail that the 49,000 training embeddings keep.
\item \textbf{Where the unseen classes land.} 90\% of the novel images fall in 71 to 97 cells, while 90\% of the normal ones need 170 to 174.
\end{itemize}

\section{Leakage Control}

\begin{itemize}
  \item The Phase 1 network and the Phase 2 map are frozen and checked by SHA-256. The SOM scaler, the Mahalanobis means and covariance, the nearest-neighbour reference, the classical detectors and the QE threshold are all fitted on the Fashion-MNIST training embeddings.
  \item The SVHN images are used only to compute scores; nothing is fitted or chosen on them, and their digit labels are not used.
  \item The scores, measures, statements and verdict rules were written in the Fashion-MNIST protocol (with amendments A1 and A2) before any Fashion-MNIST novelty score existed. The length and autoencoder diagnostics were run because the protocol says to run them if the MNIST problems reappear; they are exploratory.
  \item The held-out-class networks never see their unknown classes, and only training and validation images are used.
  \item The Fashion-MNIST test split is not read in Phase 3. Every run records the SHA-256 of its inputs, and no finished run is overwritten.
  \item \textbf{Declared:} the protocol was written after the MNIST results were known. It copies MNIST's design rather than adapting it.
\end{itemize}

\section{Limitations}

\begin{itemize}[leftmargin=1.4em, itemsep=2pt]
  \item One network and one map for the SVHN test, so the intervals show only which images were sampled. The SOM seed alone moves the QE AUROC by about 0.1.
  \item Only one outside set. SVHN is far novelty and differs mainly in embedding length, so it says little about images that look like clothing; near novelty is covered only by unseen Fashion-MNIST classes.
  \item The held-out-class splits share classes (Sandal is unseen in all five), so they are not independent.
  \item The comparisons with nine methods and the length checks are exploratory.
  \item The figures are kept in DVC; the ones shown here were copied from their run folders.
\end{itemize}

\section{Conclusion and Next Phase}

On Fashion-MNIST the SOM detects both kinds of novelty better than chance, but its strengths differ from MNIST. Against far novelty (SVHN) it is clearly behind the network's own energy and confidence scores, and its fixed alarm does not work, because SVHN's short embeddings make it look like a typical bag to the map. Against near novelty (unseen clothing classes) it beats max-softmax in every split, which it did not do on MNIST. Plain nearest-neighbour distance on the same embeddings beats the SOM in both tests, so the map's added value is again the picture it gives of \emph{where} novel images land. Phase 4 uses the map's hotspots to guide retraining and repeats these novelty tests on the retrained networks.

\clearpage
\printbibliography[title={References}]

\end{document}
